\makeatletter

\geometry{
  left=2.5cm,
  right=2.5cm,
  top=2.5cm,
  bottom=2.5cm,
  bindingoffset=0cm,
  includehead=false,
  includefoot=false,
  headheight=18pt,
  headsep=12pt,
  footskip=1.6cm
}

\KOMAoptions{
  headsepline=0.4pt
}

% --- Fonts ---------------------------------------------------------------
% Use internal helpers for detection, then define UPC-specific font families.
\thesis@setlatinfont
\thesis@setCJKfonts

\newCJKfontfamily\upcsong[
  AutoFakeBold,
  ItalicFont=\thesis@CJKkai,
  BoldItalicFont=\thesis@CJKhei
]{\thesis@CJKsong}
\newCJKfontfamily\upchei{\thesis@CJKhei}[AutoFakeBold]
\newCJKfontfamily\upckai{\thesis@CJKkai}[AutoFakeBold]
\newCJKfontfamily\upcfang{\thesis@CJKfang}[AutoFakeBold]
\newfontfamily\upctimes[
  BoldFont=\thesis@latinfont@bold,
  ItalicFont=\thesis@latinfont@italic,
  BoldItalicFont=\thesis@latinfont@bolditalic
]{\thesis@latinfont@name}

% --- Chinese font sizes --------------------------------------------------
\newcommand*\xiaosanhao{\fontsize{15pt}{22.5pt}\selectfont}
\newcommand*\xiaochuhao{\fontsize{36pt}{54pt}\selectfont}
\newcommand*\xiaoerhao{\fontsize{18pt}{27pt}\selectfont}
\newcommand*\sanhao{\fontsize{16pt}{24pt}\selectfont}
\newcommand*\sihao{\fontsize{14pt}{21pt}\selectfont}
\newcommand*\xiaosihao{\fontsize{12pt}{23.4pt}\selectfont}
\newcommand*\wuhao{\fontsize{10.5pt}{15.75pt}\selectfont}
\newlength\upctitlepagespace
\setlength\upctitlepagespace{4\baselineskip}
\newlength\upcsectionspace
\setlength\upcsectionspace{9pt}

% --- Metadata ------------------------------------------------------------
% \studentnumber 在 internal/titlepage.tex 中默认会加方括号，
% UPC 格式不需要括号，因此在这里用 \renewcommand* 覆盖。
\renewcommand*\studentnumber[1]{\def\@studentnumber{#1}}
\def\@studentnumber{}

% UPC 使用 \college 而非 \institute，故不向后者的宏传递。
\newcommand*\college[1]{\def\@college{#1}}
\def\@college{}

\newcommand*\major[1]{\def\@major{#1}}
\def\@major{}

\newcommand*\advisor[1]{\def\@advisor{#1}}
\def\@advisor{}

\newcommand*\upcsubtitle[1]{\def\@upcsubtitle{#1}}
\def\@upcsubtitle{}

\newcommand*\upcenglishtitle[1]{\def\@upcenglishtitle{#1}}
\def\@upcenglishtitle{The design and implementation of the linear form}

\newcommand*\upcenglishsubtitle[1]{\def\@upcenglishsubtitle{#1}}
\def\@upcenglishsubtitle{}

\newcommand*\upckeywords[1]{\def\@upckeywords{#1}}
\def\@upckeywords{数据结构；面向对象；可视化；算法}

\newcommand*\upcenkeywords[1]{\def\@upcenkeywords{#1}}
\def\@upcenkeywords{data structure；object-oriented；visualization；algorithm}

\newcommand*\UPCCNTitle{\@title}
\newcommand*\UPCCNSubtitle{\@upcsubtitle}
\newcommand*\UPCEnglishTitle{\@upcenglishtitle}
\newcommand*\UPCEnglishSubtitle{\@upcenglishsubtitle}
\newcommand*\UPCKeywordsCN{\@upckeywords}
\newcommand*\UPCKeywordsEN{\@upcenkeywords}

% --- Page styles ---------------------------------------------------------
\newcommand*\UPCFixedHeaderText{中国石油大学（华东）本科毕业设计 (论文)}
\newcommand*\upcheaderfont{\upckai\wuhao}
\newcommand*\upcfooterfont{\upcsong\wuhao}
\newcommand*\upcheadercontent[1]{\raisebox{-3pt}[0pt][0pt]{\upcheaderfont #1}}

\newpairofpagestyles{upcfrontmatter}{%
  \clearpairofpagestyles
  \chead{\upcheadercontent{\UPCFixedHeaderText}}
}

\newpairofpagestyles{upcfrontmatterplain}{%
  \clearpairofpagestyles
}
\ModifyLayer[contents={}]{upcfrontmatterplain.head.above.line}
\ModifyLayer[contents={}]{upcfrontmatterplain.head.below.line}
\ModifyLayer[contents={}]{plain.upcfrontmatterplain.head.above.line}
\ModifyLayer[contents={}]{plain.upcfrontmatterplain.head.below.line}

\newpairofpagestyles{upcmainmatter}{%
  \clearpairofpagestyles
  \chead{\upcheadercontent{\leftmark}}
  \cfoot{\upcfooterfont \pagemark}
}

\automark[chapter]{chapter}

\newcommand*\upcsetupfrontmatter{%
  \pagestyle{upcfrontmatter}%
  \renewcommand*\chapterpagestyle{upcfrontmatter}%
  \markboth{\UPCFixedHeaderText}{}%
}

\newcommand*\upcsetupmainmatter{%
  \pagestyle{upcmainmatter}%
  \renewcommand*\chapterpagestyle{upcmainmatter}%
}

\renewcommand*{\chaptermark}[1]{\markboth{第 \thechapter 章\enskip #1}{}}
\renewcommand*{\sectionmark}[1]{}

% --- Heading styles ------------------------------------------------------
\setkomafont{chapter}{\normalfont\upchei\sanhao}
\setkomafont{section}{\normalfont\upchei\sihao}
\setkomafont{subsection}{\normalfont\upchei\xiaosihao}
\setkomafont{subsubsection}{\normalfont\upcsong\xiaosihao}

\RedeclareSectionCommand[beforeskip=0.5\baselineskip,afterskip=1.5\baselineskip]{chapter}
\RedeclareSectionCommand[beforeskip=\upcsectionspace,afterskip=\upcsectionspace]{section}
\RedeclareSectionCommand[beforeskip=\upcsectionspace,afterskip=\upcsectionspace]{subsection}
\RedeclareSectionCommand[beforeskip=\upcsectionspace,afterskip=\upcsectionspace]{subsubsection}

% --- Figure / table / equation styles ------------------------------------
\DeclareCaptionFont{upcwuhao}{\upcsong\wuhao}
\DeclareCaptionLabelFormat{upc}{#1 #2}
\captionsetup{
  font=upcwuhao,
  labelfont=upcwuhao,
  textfont=upcwuhao,
  labelformat=upc,
  labelsep=space,
  justification=centering,
  singlelinecheck=false
}

% --- TOC styles ----------------------------------------------------------
\newcommand*\upc@toctitle{目 \quad 录}
\setkomafont{chapterentry}{\upcsong\xiaosihao}
\setkomafont{chapterentrypagenumber}{\upcsong\xiaosihao}

\newcommand*\upc@tocchapterentrynumber[1]{第 #1 章\enskip}
\newcommand*\upc@tocnumber[1]{#1\enskip}
\newcommand*\upc@tocentry[1]{\upcsong\xiaosihao #1}
\newcommand*\upc@tocpagenumber[1]{\upcsong\xiaosihao #1}

\DeclareTOCStyleEntry[
  entrynumberformat=\upc@tocchapterentrynumber,
  entryformat=\upc@tocentry,
  pagenumberformat=\upc@tocpagenumber,
  linefill=\TOCLineLeaderFill,
  raggedpagenumber=false,
  indent=0pt,
  numwidth=3.5em,
  beforeskip=0pt
]{tocline}{chapter}

\DeclareTOCStyleEntry[
  entrynumberformat=\upc@tocnumber,
  entryformat=\upc@tocentry,
  pagenumberformat=\upc@tocpagenumber,
  linefill=\TOCLineLeaderFill,
  raggedpagenumber=false,
  indent=2em,
  numwidth=2em,
  beforeskip=0pt
]{tocline}{section}

\DeclareTOCStyleEntry[
  entrynumberformat=\upc@tocnumber,
  entryformat=\upc@tocentry,
  pagenumberformat=\upc@tocpagenumber,
  linefill=\TOCLineLeaderFill,
  raggedpagenumber=false,
  indent=4em,
  numwidth=3em,
  beforeskip=0pt
]{tocline}{subsection}

% --- Utility commands ----------------------------------------------------
\newcommand*\upcfronttitle[1]{%
  \cleardoublepage
  \vspace*{\baselineskip}%
  {\centering \upchei\xiaosanhao #1\par}%
  \vspace{0.5\baselineskip}%
}

\newcommand*\upcabstractheading[1]{%
  \vspace{0.5\baselineskip}%
  {\centering #1\par}%
  \vspace{0.5\baselineskip}%
}

\newcommand*\upcspecialchapter[2]{%
  \cleardoublepage
  \upcsetupfrontmatter
  \chapter*{#1}%
  \thispagestyle{upcfrontmatter}%
  \addcontentsline{toc}{chapter}{#2}%
}

\newcommand*\upcprinttableofcontents{%
  \cleardoublepage
  \upcsetupfrontmatter
  \begingroup
  \renewcommand*\contentsname{\upc@toctitle}%
  \RedeclareSectionCommand[beforeskip=1.5\baselineskip,afterskip=0.5\baselineskip]{chapter}%
  \tableofcontents
  \endgroup
  \thispagestyle{upcfrontmatter}%
}

\newcommand*\upcappendixnumbering{%
  \setcounter{figure}{0}%
  \setcounter{table}{0}%
  \setcounter{equation}{0}%
  \renewcommand*\thefigure{A\arabic{figure}}%
  \renewcommand*\thetable{A\arabic{table}}%
  \renewcommand*\theequation{A\arabic{equation}}%
}

\newcommand*\upcbody{\upcsong\xiaosihao}
\newcommand*\upcabstractbody{\upcbody}
\newcommand*\upcbodybf{\upcsong\xiaosihao\bfseries}
\AtBeginDocument{\upcbody}
\newcommand*\upccoverline[2][7cm]{\underline{\makebox[#1][l]{\strut #2}}}

% --- Bibliography & citation styles --------------------------------------
% 参考文献列表：宋体小四，取消所有斜体映射以避免 fallback 到楷体
\AtBeginBibliography{%
  \upcsong\xiaosihao
  \renewcommand*{\mkbibemph}[1]{#1}%
  \renewcommand*{\textit}[1]{#1}%
  \renewcommand*{\emph}[1]{#1}%
}

% 引用：小四号 Times New Roman 上标
\DeclareCiteCommand{\cite}[\upcbibciteformat]
  {\bibopenbracket}
  {\usebibmacro{citeindex}\usebibmacro{cite}}
  {\bibclosebracket\multicitedelim\bibopenbracket}
  {\bibclosebracket}
\newcommand*\upcbibciteformat[1]{\textsuperscript{\upctimes\xiaosihao #1}}

\newcommand*\upccoverlinefill[2][7cm]{%
  \if\relax\detokenize{#2}\relax
    \upccoverline[#1]{}%
  \else
    \upccoverline[#1]{#2}%
  \fi
}

% ---------------------------------------------------------------------------
% Decoupled environments and commands for segments/upc
% ---------------------------------------------------------------------------

% --- Abstracts -------------------------------------------------------------
\newenvironment{upcabstractcn}{%
  \cleardoublepage
  \upcsetupfrontmatter
  \thispagestyle{upcfrontmatter}%
  \markboth{\UPCFixedHeaderText}{}%
  \vspace*{\baselineskip}%
  {\centering {\upchei\xiaoerhao \UPCCNTitle\par}}%
  \ifdefempty{\@upcsubtitle}{}{\par\raggedleft\upchei\xiaoerhao \UPCCNSubtitle\par}%
  \upcabstractheading{\upchei\sanhao 摘\quad 要}%
  \upcabstractbody
}{%
  \par
  \vspace{0.5\baselineskip}%
  {\upcabstractbody\textbf{关键词：}\UPCKeywordsCN\par}%
}

\newenvironment{upcabstracten}{%
  \cleardoublepage
  \upcsetupfrontmatter
  \thispagestyle{upcfrontmatter}%
  \markboth{\UPCFixedHeaderText}{}%
  \vspace*{\baselineskip}%
  {\centering {\upctimes\xiaoerhao\bfseries \UPCEnglishTitle\par}}%
  \ifdefempty{\@upcenglishsubtitle}{}{\par\raggedleft\upctimes\xiaoerhao\bfseries \UPCEnglishSubtitle\par}%
  \upcabstractheading{\upctimes\sanhao\bfseries Abstract}%
  \upcabstractbody
}{%
  \par
  \vspace{0.5\baselineskip}%
  {\upcabstractbody\textbf{Keywords: }\UPCKeywordsEN\par}%
}

% --- Originality & License -------------------------------------------------
\newcommand*\UPCOriginalityTitle{学位论文原创性声明}
\newcommand*\UPCOriginalityText{%
  本人所提交的学位论文《\@title》，是在导师的指导下，独立进行研究工作所取得的原创性成果。除文中已经注明引用的内容外，本论文不包含任何其他个人或集体已经发表或撰写过的研究成果。对本文的研究做出重要贡献的个人和集体，均已在文中标明。

  本声明的法律后果由本人承担。%
}
\newcommand*\UPCAuthorSignatureLabel{论文作者（签名）：}
\newcommand*\UPCAdvisorConfirmLabel{指导教师确认（签名）：}
\newcommand*\UPCLicenseTitle{学位论文版权使用授权书}
\newcommand*\UPCLicenseText{%
  本学位论文作者完全了解中国石油大学（华东）有权保留并向国家有关部门或机构送交学位论文的复印件和磁盘，允许论文被查阅和借阅。本人授权中国石油大学（华东）可以将学位论文的全部或部分内容编入有关数据库进行检索，可以采用影印、缩印或其它复制手段保存、汇编学位论文。%
}
\newcommand*\UPCAdvisorSignatureLabel{指导教师（签名）：}
\newcommand*\UPCDatePlaceholder{\quad 年 \quad 月 \quad 日}

\newcommand*\upc@signatureblock[2]{%
  \noindent
  \begin{minipage}[t]{0.48\textwidth}%
    {\upcbody\centering #1\par}%
  \end{minipage}\hfill
  \begin{minipage}[t]{0.48\textwidth}%
    {\upcbody\centering #2\par}%
  \end{minipage}%
  \noindent
  \begin{minipage}[t]{0.48\textwidth}%
    {\upcbody\raggedleft \UPCDatePlaceholder\par}%
  \end{minipage}\hfill
  \begin{minipage}[t]{0.48\textwidth}%
    {\upcbody\raggedleft \UPCDatePlaceholder\par}%
  \end{minipage}%
}

\newcommand*\upcprintoriginality{%
  \cleardoublepage
  \thispagestyle{upcfrontmatterplain}%
  \vspace{\baselineskip}%
  \begin{center}
    {\upchei\xiaosanhao \UPCOriginalityTitle\par}%
  \end{center}%
  \vspace{\baselineskip}%
  {\upcbody \UPCOriginalityText\par}%
  \vspace{4\baselineskip}%
  \upc@signatureblock{\UPCAuthorSignatureLabel}{\UPCAdvisorConfirmLabel}%
}

\newcommand*\upcprintlicense{%
  \vspace{2\baselineskip}%
  \begin{center}
    {\upchei\xiaosanhao \UPCLicenseTitle\par}%
  \end{center}%
  \vspace{\baselineskip}%
  {\upcbody \UPCLicenseText\par}%
  \vspace{4\baselineskip}%
  \upc@signatureblock{\UPCAuthorSignatureLabel}{\UPCAdvisorSignatureLabel}%
}

% --- Appendix --------------------------------------------------------------
\newcommand*\upcappendixprefix{%
  \upcspecialchapter{附　录}{附　录}%
  \upcappendixnumbering
}

\newcommand*\upcappendixsection[1]{%
  \section*{#1}%
  \addcontentsline{toc}{section}{#1}%
}

% --- Table helpers (three-line table) --------------------------------------
% UPC table spec:
%   - 5号宋体 (\upcsong\wuhao)
%   - 三线表 (\toprule / \midrule / \bottomrule)
%   - 表号按一级标题编排 (already: \thechapter-\arabic{table})
%   - 表头单位换行（与中文名不在同一行）
%   - 跨页表格自动输出“续表××”
%   - 表注写在底线下左侧

% 统一表格字体
\newcommand*\upctablefont{\upcsong\wuhao}

% 表头单元格：中文名 + 换行 + 单位
% 用法：\upchdr{试验编号}{/（个）}
\newcommand*\upchdr[2]{%
  \begin{tabular}[c]{@{}c@{}}\strut #1\\#2\strut\end{tabular}%
}

% 普通三线表环境
% #1[可选] = 浮动位置 (默认 htbp)
% #2 = 表题
% #3 = label
\newenvironment{upctable}[3][htbp]{%
  \begin{table}[#1]
  \centering
  \upctablefont
  \setlength{\tabcolsep}{3.5pt}%
  \renewcommand{\arraystretch}{1.2}%
  \caption{#2}\label{#3}%
}{%
  \end{table}%
}

% 跨页三线表环境
% 封装了字体、列距、行距和 caption，用户按照标准 longtable 语法填写头部与数据即可。
% #1 = 列格式定义 (e.g. {cc})
% #2 = 表题
% #3 = label
\newenvironment{upclongtable}[3]{%
  \upctablefont
  \setlength{\tabcolsep}{3.5pt}%
  \renewcommand{\arraystretch}{1.2}%
  \begin{longtable}{#1}%
  \caption{#2}\label{#3}\\%
}{%
  \end{longtable}%
}

% 续表标题行（放在 \endfirsthead 之后、续页表头之前）
% #1 = 总列数（用于 \multicolumn）
\newcommand*\upclongtablecont[1]{%
  \multicolumn{#1}{r}{续表~\thetable}\\%
}

% 表注：放在表格底线下左侧
% #1[可选] = 合并列数（默认 1）
% #2 = 注释内容
% 用法：跨页表放在数据最后、\bottomrule 之前；普通表放在 \end{tabular} 之前
\newcommand*\upctablenote[2][1]{%
  \\[-0.6em]%
  \multicolumn{#1}{@{}l}{\upctablefont 注：#2}%
}

\makeatother
